\section{AI/ML Forecasting}
\label{sec:ai_ml_forecasting}

\subsection{Forecasting Objective \& Approach}
\label{subsec:forecasting_objective}

The forecasting engine estimates hourly renewable power generation $\hat{y}_t$ across 24-hour, 48-hour, and 72-hour horizons. For the prototype implementation, GridSense AI adopts **XGBoost (eXtreme Gradient Boosting)**, leveraging its proven efficiency on tabular time-series features, robust non-linear interaction modeling, and rapid training cycles.

\textbf{Feature Engineering Pipeline:}
\begin{itemize}[leftmargin=1.5em, itemsep=0.2em]
    \item \textbf{Weather Variables:} Global Horizontal Irradiance (GHI), Direct Normal Irradiance (DNI), ambient temperature, cloud cover fraction, wind velocity at hub height, and air density.
    \item \textbf{Historical \& Temporal Lags:} Autoregressive power lags ($y_{t-1}, y_{t-24}, y_{t-48}$), 6-hour rolling statistics, hour-of-day solar zenith angles, and seasonal day-of-year encodings.
    \item \textbf{Site Constraints:} Installed nameplate capacity ($P_{\text{rated}}$), inverter thresholds, and tilt/hub height metadata.
\end{itemize}

\subsection{Forecast Uncertainty \& Risk Quantification}
\label{subsec:uncertainty_risk}

Rather than outputting a single deterministic curve, GridSense AI computes quantile prediction intervals ($p_{10}$, $p_{50}$, $p_{90}$) to quantify forecast dispersion (\cref{fig:forecast_uncertainty_chart}):
\begin{center}
    \texttt{Expected Generation } $= 80\text{ MW} \quad \big\vert \quad \texttt{Prediction Range } [p_{10}, p_{90}] = [72\text{ MW}, 88\text{ MW}]$
\end{center}

\begin{figure}[htbp]
    \centering
    \begin{tikzpicture}[xscale=0.80, yscale=0.62]
        % Background grid
        \draw[very thin, color=borderGray!50, step=1cm] (0,0) grid (12,5);
        \draw[->, thick, color=darkSlate] (0,0) -- (12.3,0) node[right] {\footnotesize Time};
        \draw[->, thick, color=darkSlate] (0,0) -- (0,5.3) node[above] {\footnotesize Generation (MW)};
        
        % Ticks
        \foreach \x/\label in {0/00:00, 3/06:00, 6/12:00, 9/18:00, 12/24:00}
            \draw (\x,0.1) -- (\x,-0.1) node[below] {\scriptsize \label};
        \foreach \y/\label in {1/20, 2/40, 3/60, 4/80, 5/100}
            \draw (0.1,\y) -- (-0.1,\y) node[left] {\scriptsize \label};

        % Shaded Uncertainty Envelope (p10 to p90 band)
        \fill[accentTeal!18] 
            (0,0) -- 
            plot [smooth] coordinates {(2.5,0) (4,1.8) (6,4.8) (7.5,3.6) (9,1.6) (10,0)} --
            plot [smooth] coordinates {(10,0) (9,0.4) (7.5,1.8) (6,3.6) (4,0.6) (2.5,0)} -- cycle;

        % Net Demand Curve
        \draw[dashed, thick, color=darkSlate!70] plot [smooth] coordinates {(0,2.2) (3,1.8) (6,3.2) (9,4.5) (12,2.8)};
        \node[color=darkSlate!80, font=\scriptsize] at (10.4,4.7) {Scheduled Demand};

        % Expected Forecast Line (p50)
        \draw[line width=1.3pt, color=accentTeal] plot [smooth] coordinates {(0,0) (2.5,0) (4,1.2) (6,4.2) (7.5,2.7) (9,0.9) (10,0)};
        \node[color=accentTeal, font=\scriptsize\bfseries] at (6,4.5) {Forecast ($p_{50}$)};

        % Highlighted Risk / Deficit Zone
        \draw[fill=accentRose!20, draw=accentRose, dashed, line width=0.8pt] (7.5,0.2) rectangle (9.8,4.6);
        \node[color=accentRose, font=\scriptsize\bfseries, align=center] at (8.65,2.5) {High Shortfall\\Risk Window};
    \end{tikzpicture}
    \caption{Conceptual 24-hour renewable generation forecast with probabilistic prediction intervals ($p_{10}$--$p_{90}$) and highlighted shortfall risk window.}
    \label{fig:forecast_uncertainty_chart}
\end{figure}

Uncertainty is translated into operational risk metrics by evaluating the dispersion width ($p_{90} - p_{10}$) alongside supply-demand margins ($P_{\text{gen}} < P_{\text{demand}}$).

\subsection{Explainability \& Model Evaluation}
\label{subsec:eval_metrics}

The model's predictions are explained in natural language (e.g., \emph{``Generation drops 40\% at 15:00 due to rapid cloud formation''}). Offline validation employs standard regression metrics (MAE, RMSE, MAPE) with **time-aware expanding-window cross-validation** to guarantee that future meteorological information never leaks into historical training sets.
